% =====================================================================
%  TITRE PRÉLIMINAIRE : DISPOSITIONS GÉNÉRALES
% =====================================================================
\titrepartie{Titre préliminaire}{Dispositions générales}

% ---------------------------------------------------------------------
\article{Objet, champ d'application et portée}

\chapeau{Le présent article détermine à qui s'applique le règlement, ce qu'il contient et quelle est sa valeur au regard des autres textes du département.}

\cl[objet]{Le présent règlement fixe les règles d'organisation et de fonctionnement du Los Santos Police Department, les obligations déontologiques et professionnelles de ses personnels, les procédures opérationnelles de référence ainsi que le régime disciplinaire applicable en cas de manquement.}

\cl[champ]{Il s'applique à l'ensemble des personnels du département, sans exception : agents assermentés de tout grade, personnels administratifs, membres des divisions et unités spécialisées, agents en période probatoire et agents détachés auprès d'un autre service.}

\cl{L'engagement au sein du département vaut acceptation sans réserve du présent règlement. Nul ne peut invoquer son ignorance, son ancienneté, son grade ou l'usage en vigueur dans son poste d'affectation pour justifier un manquement.}

\cl[primaute]{Le présent règlement prime sur toute pratique orale, coutume de poste ou consigne informelle contraire. Les notes de service, protocoles de division et instructions permanentes le complètent sans jamais le contredire ; en cas de contradiction, la disposition du présent règlement prévaut jusqu'à sa modification formelle par l'État-Major.}

\cl{Le règlement ne se substitue ni au Code pénal, ni au Code de la route, ni au Code de procédure applicable dans l'État de San Andreas. Lorsqu'une disposition du présent règlement est plus exigeante que la loi, c'est la disposition la plus exigeante qui s'applique à l'agent.}

\cl{Chaque agent est responsable de détenir la version en vigueur du règlement et de prendre connaissance de ses mises à jour dès leur diffusion. La consultation du registre des versions figurant en tête du présent document suffit à cette vérification.}

\begin{note}[Citer une disposition]
Une disposition se cite par son numéro d'alinéa entre crochets, éventuellement suivi de son intitulé d'article. Exemple : \emph{« manquement à l'alinéa \refcl{primaute}, article “Objet, champ d'application et portée” »}. Toute demande de sanction, tout recours et tout rapport disciplinaire doit viser au moins une disposition précise.
\end{note}

% ---------------------------------------------------------------------
\article{Définitions}

\chapeau{Les termes ci-après ont, dans l'ensemble du présent règlement, le sens qui leur est donné au présent article.}

\begin{definitions}
\dterme{Agent} Toute personne employée par le département, quel que soit son grade ou sa fonction. Les termes \emph{officier}, \emph{policier}, \emph{personnel}, \emph{employé} et \emph{membre} lui sont équivalents dans l'ensemble du texte.
\dterme{Agent assermenté} Agent ayant prêté le serment d'assermentation, doté d'une plaque et habilité à exercer les prérogatives de police.
\dterme{État-Major} Le Chief et les Commanders. Autorité supérieure du département.
\dterme{Command Staff} Les Captains et les Lieutenants. Encadrement supérieur des postes et des divisions.
\dterme{Supervision} Les Sergeants I et II, ainsi que, pour les actes qui leur sont expressément ouverts, les Senior Lead Officers.
\dterme{Superviseur} Tout agent titulaire d'un grade de supervision ou supérieur, exerçant effectivement l'autorité sur une intervention ou sur une unité.
\dterme{Watch Commander} Superviseur de référence désigné pour une garde, responsable de l'organisation et de la supervision des unités en service.
\dterme{Operation Lead} Agent qui exerce le commandement d'une intervention en cours. Par défaut, le superviseur présent le plus élevé en grade ; à défaut de superviseur, l'agent primo-intervenant jusqu'à relève.
\dterme{Unité} Équipage constitué d'un ou plusieurs agents disposant d'un indicatif radio propre.
\dterme{Unité primaire} Unité qui conduit une intervention, une poursuite ou une procédure et qui en assume la responsabilité radio.
\dterme{Garde} Période de service continue d'un agent, du début à la fin de sa prise de service.
\dterme{Suspicion raisonnable} Ensemble de faits précis et exprimables, et des déductions rationnelles qui peuvent en être tirées, permettant de soupçonner qu'une infraction a été, est ou va être commise. Elle autorise le contrôle et la palpation de sécurité.
\dterme{Cause probable} Ensemble de faits qui conduiraient une personne raisonnable à croire qu'une infraction a été commise et que la personne visée en est l'auteur. Elle est requise pour procéder à une arrestation, à une fouille ou à une saisie.
\dterme{Usage de la force} Tout recours à la contrainte physique, à une arme intermédiaire ou à une arme à feu à l'encontre d'une personne.
\dterme{Force létale} Force susceptible d'entraîner la mort ou des blessures graves.
\dterme{Note de service} Décision écrite de l'État-Major ou du Command Staff, motivée, diffusée à l'ensemble des personnels concernés, et venant préciser l'application du présent règlement.
\dterme{Division majeure} Division dont l'intégration est exclusive : un agent ne peut appartenir qu'à une seule division majeure.
\dterme{Division mineure} Division dont l'intégration est cumulable avec d'autres divisions mineures.
\dterme{First Lincoln} Évaluation finale de la période probatoire, à l'issue de laquelle l'agent est autorisé à patrouiller sans encadrement direct.
\end{definitions}

\begin{note}[Rédaction inclusive]
Les fonctions, grades et qualités mentionnés au présent règlement s'entendent indifféremment au masculin et au féminin. L'emploi du masculin générique répond à une exigence de lisibilité et n'emporte aucune distinction de traitement.
\end{note}
